\section{Contoh}
\label{sec:contoh-bab11}

Dua contoh berikut memperlihatkan manipulasi bit pada program utuh.
Contoh pertama memakai \asm{SHR} dan \asm{TEST} untuk memeriksa
paritas. Contoh kedua memakai \asm{XOR} untuk enkripsi dan dekripsi
sederhana.

\subsection{Contoh 1: Memeriksa Paritas}

Program pada Kode~\ref{lst:paritas} menghitung jumlah bit 1 pada
sebuah byte, lalu menentukan paritasnya.

\lstinputlisting[
    language=MASM,
    caption={Program \file{paritas.asm} memeriksa paritas byte.},
    label={lst:paritas}
]{\CodeDir/paritas.asm}

Bagian penting program:

\begin{itemize}
    \item \lstasm|MOV CX, 8| menyiapkan delapan kali pergeseran.
    \item \lstasm|SHR AL, 1| menggeser bit terkanan ke \flag{CF}.
    \item \lstasm|ADC BL, 0| menambahkan \flag{CF} ke penghitung,
          sehingga \reg{BL} berisi jumlah bit 1.
    \item \lstasm|TEST BL, 1| memeriksa bit terendah jumlah bit.
          Bila bit itu 1, jumlahnya ganjil; bila 0, jumlahnya genap.
\end{itemize}

Byte uji \bin{10110110} memiliki lima bit 1, sehingga paritasnya
ganjil. Keluaran program tampak pada
Kode~\ref{lst:keluaran-paritas}.

\begin{lstlisting}[
    style=sessiongaya,
    caption={Keluaran program \file{paritas.asm}.},
    label={lst:keluaran-paritas}
]
Byte 10110110b berparitas: GANJIL
\end{lstlisting}

\subsection{Contoh 2: XOR Cipher}

Program pada Kode~\ref{lst:xor-cipher} memakai \asm{XOR} untuk
mengenkripsi dan mendekripsi string. Sifat $A \oplus K \oplus K = A$
membuat proses dapat dibalik dengan kunci yang sama.

\lstinputlisting[
    language=MASM,
    caption={Program \file{xor-cipher.asm} mengenkripsi dan
    mendekripsi string dengan XOR.},
    label={lst:xor-cipher}
]{\CodeDir/xor-cipher.asm}

Bagian penting program:

\begin{itemize}
    \item \lstasm|MOV AL, [SI]| mengambil satu byte dari string asli.
    \item \lstasm|XOR AL, kunci| membalik bit sesuai kunci.
    \item \lstasm|MOV [DI], AL| menyimpan hasil ke penyangga.
    \item Pemeriksaan \lstasm|CMP AL, '$'| menghentikan perulangan
          pada akhir string.
    \item Dekripsi memakai \asm{XOR} yang sama pada hasil enkripsi.
\end{itemize}

Dengan kunci \hex{20}, huruf besar berubah menjadi huruf kecil
karena bit ke-5 dibalik. Keluaran program tampak pada
Kode~\ref{lst:keluaran-xor}.

\begin{lstlisting}[
    style=sessiongaya,
    caption={Keluaran program \file{xor-cipher.asm}.},
    label={lst:keluaran-xor}
]
Teks asli    : BELAJARASM
Terenkripsi  : belajarasm
Terdekripsi  : BELAJARASM
\end{lstlisting}

\begin{infobox}
\vspace{1.1em}
\eng{XOR cipher} pada contoh ini hanya demonstrasi sifat \asm{XOR},
bukan algoritma kriptografi yang aman. Kunci satu byte dan pola
sederhana mudah ditebak; pembahasan keamanan berada di luar cakupan
buku ini.
\end{infobox}
